\subsection{Classifier Replay by Year}
\label{sec:results:replay}

Table~\ref{tab:replay} and Figure~\ref{fig:replay} give the detection rate of the deterministic rule by year for each series next to the published o4-mini rates (per-year values for all windows of each series are in Table~S8 of the Supplementary Materials), and Figure~S5 of the Supplementary Materials shows how far the replay depends on the thresholds.
% src: stored_e3_intervals.csv source=rule (moving_block block_len=60 ci_lo_pct 3.78, ci_hi_pct 24.10; wilson 9.40 to 15.08); stored_e3_ess.csv source=rule (n_episodes 7, ess_acf 31.7, ess_batch_30 29.4, ess_batch_sqrt_n 41.8); stats.yaml rates arm=r1_o4mini periods 2024, 2025 (wilson_ci_pct); report.yaml headline.llm_vs_classifier_all_years arm=r1_o4mini (agree_pct, kappa, n_compared); detection_by_year.csv (classifier rows by series_key and year; arm r1_o4mini); stats.yaml rates (Wilson, bootstrap_quarter, nonoverlap_30d) and contrasts series_contrast period 2024-2025 and period_shift
On the recovered inputs the rule reproduces the published labels in 1{,}412 of 1{,}412 parsed windows (100.0\%; Cohen's kappa 1.000), so the published labels follow from the published inputs and the published criteria. On the published inputs the rule detects every window ending in 2024 and 2025, and so do the published o4-mini answers (Table~\ref{tab:replay}; Wilson 95\% lower bound 98.4\% in 2024 and 98.5\% in 2025). The call-only series and the archive give the same saturation. On the complete series the rule detects 10.3\% in 2024 and 13.6\% in 2025 (Table~\ref{tab:replay}).

Pooled over 2024 and 2025 the detection rate is 100.0\% on the published inputs (501 of 501 windows) and 12.0\% on the complete series (60 of 502). Because the windows overlap at stride one, the interval that we report is the 60-window moving-block bootstrap interval, 3.8--24.1\%; the Wilson interval, 9.4--15.1\%, treats the 502 windows as independent and is too narrow. The 60 detections form 7 episodes, and the effective sample size of the 502 windows is 29 to 42 by three estimators (Section~S9 of the Supplementary Materials); the non-overlapping subsample holds 3 detections in 17 windows (5.9\% to 25.0\% over the 30 starts). The two inputs differ by 88.0 percentage points on the same windows (calendar-quarter bootstrap interval 80.4--94.9, 8 blocks). This is a difference between two inputs on the same windows, not between two samples drawn from a population, so we give no p-value for it.

% src: detection_by_year.csv label_source=classifier symbol=SPY scope=year series_key=jan2026_phase4a|new_full years 2022, 2024, 2025, 2026 (n_detect_positive, n_detect_negative); condensed/windows.csv.gz series_key=call_only (clf_positive_days 30 and clf_sign_flips 0 in all 1,672 windows; the 21 windows not detected are low conviction, ends 2020-04-02 to 2020-05-01) and series_key=new_full symbol=SPY (window_end of the detected persistent_negative windows of 2025: 2025-03-18 to 2025-04-29)
The detections also differ in sign. On the published inputs all 501 detected windows of 2024 and 2025 are persistent positive-GEX windows. The call leg alone is positive by construction (Equation~(\ref{eq:gex})): each of its 1{,}672 windows has 30 positive days and no sign change, so the magnitude criterion alone decides, and it fails only in the 21 windows that end from 2020-04-02 to 2020-05-01. On the complete series the 60 detections of 2024 and 2025 are 30 positive windows (26 ending in 2024 and 4 in 2025) and 30 negative windows, all ending between 2025-03-18 and 2025-04-29. The detections of 2022 (102 windows) and of 2026 (22 windows; Section~\ref{sec:results:post}) are all negative. The published inputs therefore implied a positive regime in every window of 2024 and 2025, which the complete series does not show.

The period contrast changes sign. On the published inputs the rule detects 17.0\% of the windows of 2020 to 2023 and 100.0\% of those of 2024 and 2025, a shift of 83.0 percentage points ($\varphi=0.789$; calendar-quarter bootstrap $p\le5.0\times10^{-4}$, the floor of 2{,}000 draws, which treats the quarters as independent; Fisher exact test on the non-overlapping subsample $p=2.7\times10^{-9}$, which treats those windows as independent). On the complete series the rates are 22.2\% and 12.0\%, a shift of $-$10.2 percentage points (60-window moving-block bootstrap interval $-$20.9 to 8.5; $\varphi=-0.124$; the calendar-quarter bootstrap, $-$23.2 to 2.0 with $p=0.108$, and the Fisher test on the subsample, $p=0.728$, are secondary). At the nominal rule the complete series therefore shows no evidence of a rise after 2023. The interval reaches a rise of 8.5 points, above a margin of 5 points, so the data do not rule out a rise of 5 points or more (other thresholds, baselines and shared windows: Section~S9 of the Supplementary Materials); this is a failure to find a rise, not a demonstration of equality. In 2020 to 2023 the two series are of the same order (Table~\ref{tab:replay}).
% src: stored_e4_shift_margin.csv threshold_bn=5.0 block_len=60 circular=False (window_set=all baseline=2020-2023: shift_pp -10.24, ci_lo_pp -20.89, ci_hi_pp 8.47; baseline=2008-2019: 1.78, -7.69, 14.62; window_set=shared baseline=2020-2023: -7.92, -20.95, 8.62), stored_e4_summary.csv window_set=all baseline=2020-2023 block_len=60 circular=False (n_thresholds 12, n_upper_below_margin 1), margin_pp 5.0; results/stats.yaml headline_contrasts name=period_shift series_key=new_full symbol=SPY (diff_pp -10.236863, bootstrap_diff_ci_pp [-23.167616, 2.001953], p_bootstrap 0.108, p_nonoverlap 0.727806); detection_by_year.csv series_key=new_full year=2020 (56 of 252) and condensed/windows.csv.gz series_key=new_full symbol=SPY window_end 2020-01-02 to 2020-02-12 (29 windows, all clf_should_detect) and window_end >= 2020-02-13 in 2020 (27 of 223)
The 2020 value of the complete series needs one qualification. The complete series reaches back before the first window of the published inputs, which ends on 2020-02-13: it holds 252 windows that end in 2020, and the 29 that end from 2020-01-02 to 2020-02-12 are all detected. The pooled 22.2\% of 2020 to 2023 contains them. On the 223 windows that the two series share, the complete series detects 12.1\% in 2020 (27 of 223), the same as the published inputs (Table~\ref{tab:replay} and Figure~\ref{fig:replay} use the shared windows; Table~S8 gives both sets).

% src: figures/output/fig12_threshold_sensitivity.json extra.rates (series_key jan2026_phase4a and new_full, rate_pct over 45 configurations) and extra.gap_pp_published_minus_rebuilt
Across the 45 threshold sets of Figure~S5 the rate on the published inputs stays at 100.0\%, whereas the rate on the complete series runs from 4.0\% to 29.3\%. The gap between the two is at least 70.7 and at most 96.0 percentage points in every configuration, so the difference between the published and the complete inputs does not depend on the exact thresholds.

\begin{table}[!htbp]
\centering
\caption{Detection rate of the deterministic classifier by year, SPY, in percent, next to the published o4-mini rate. Each cell gives the rate and, in parentheses, the windows detected out of the windows classified, counted on the windows of the series that end in the year. In 2020 every series is counted on the 223 windows that end from 2020-02-13, the span of the published inputs (the values on all windows ending in 2020 are in Table~S8); the January and archive series lack the end dates 2023-07-12 and 2024-03-11 and so hold one window fewer in 2023 and 2024. The published rate covers the 1{,}412 answers of the published result (213 of the 223 windows of 2020). The January columns of 2021 to 2023 include the call-only days of Table~\ref{tab:gaps}.}
\label{tab:replay}
\scriptsize
\setlength{\tabcolsep}{3.5pt}
\begin{tabular}{@{}lrrrrrr@{}}
\toprule
Year & Published & \ident{jan2026_phase4a} & \ident{call_only} & \ident{old_bs} & \ident{new_ex0dte} & \ident{new_full} \\
\midrule
% src: detection_by_year.csv year=2020: label_source=llm arm=r1_o4mini (26/213), classifier series_key=jan2026_phase4a (27/223) and call_only (202/223); old_bs, new_ex0dte and new_full restricted to window_end >= 2020-02-13 from condensed/windows.csv.gz (symbol SPY, clf_should_detect: 20/223, 6/223, 27/223); the all-window values of 2020 (42/252, 27/252, 56/252) are in Table S8 (tab:app:replay)
2020 & 12.2 (26/213) & 12.1 (27/223) & 90.6 (202/223) & 9.0 (20/223) & 2.7 (6/223) & 12.1 (27/223) \\
% src: detection_by_year.csv year=2021: label_source=llm arm=r1_o4mini; label_source=classifier series_key=jan2026_phase4a, call_only, old_bs, new_ex0dte, new_full (n_detect, n, rate_pct); the denominator of the new_full cell is the number of windows of the year
2021 & 3.7 (9/241) & 4.0 (10/252) & 100.0 (252/252) & 3.6 (9/252) & 0.0 (0/252) & 3.6 (9/252) \\
% src: detection_by_year.csv year=2022: label_source=llm arm=r1_o4mini; label_source=classifier series_key=jan2026_phase4a, call_only, old_bs, new_ex0dte, new_full (n_detect, n, rate_pct); the denominator of the new_full cell is the number of windows of the year
2022 & 32.4 (79/244) & 32.3 (81/251) & 100.0 (251/251) & 26.3 (66/251) & 36.7 (92/251) & 40.6 (102/251) \\
% src: detection_by_year.csv year=2023: label_source=llm arm=r1_o4mini; label_source=classifier series_key=jan2026_phase4a, call_only, old_bs, new_ex0dte, new_full (n_detect, n, rate_pct); the denominator of the new_full cell is the number of windows of the year
2023 & 20.2 (46/228) & 19.3 (48/249) & 100.0 (250/250) & 22.5 (56/249) & 25.2 (63/250) & 22.4 (56/250) \\
% src: detection_by_year.csv year=2024: label_source=llm arm=r1_o4mini; label_source=classifier series_key=jan2026_phase4a, call_only, old_bs, new_ex0dte, new_full (n_detect, n, rate_pct); the denominator of the new_full cell is the number of windows of the year
2024 & 100.0 (241/241) & 100.0 (251/251) & 100.0 (252/252) & 100.0 (251/251) & 13.1 (33/252) & 10.3 (26/252) \\
% src: detection_by_year.csv year=2025: label_source=llm arm=r1_o4mini; label_source=classifier series_key=jan2026_phase4a, call_only, old_bs, new_ex0dte, new_full (n_detect, n, rate_pct); the denominator of the new_full cell is the number of windows of the year
2025 & 100.0 (245/245) & 100.0 (250/250) & 100.0 (250/250) & 100.0 (250/250) & 16.4 (41/250) & 13.6 (34/250) \\
\bottomrule
\end{tabular}
\end{table}

\begin{figure}[!htbp]
\centering
\resultfigure{\figReplay}
\caption{Detection rate by year, SPY, 2020--2025: the published o4-mini labels and the deterministic rule replayed on the published inputs (\ident{jan2026_phase4a}), the call leg only (\ident{call_only}), the pre-rebuild archive (\ident{old_bs}) and the complete series without and with same-day expiries (\ident{new_ex0dte}, \ident{new_full}), with 95\% Wilson intervals. \textbf{Read this figure as:} the published labels and the rule on the published inputs agree window by window. In 2024 and 2025 the bars of the published inputs, the call-only series and the archive stand at 100\%, and the two bars of the complete series stand at 13.1\% and 10.3\% in 2024 and 16.4\% and 13.6\% in 2025. In 2020 the bars of the archive and of the complete series count the 223 windows that end from 2020-02-13, the span of the published inputs; on all 252 windows the complete series gives 22.2\% (Table~S8). The Wilson intervals are optimistic because the overlapping windows of a year are treated as independent.}
\label{fig:replay}
\end{figure}


